\chapter{Genealogical Path Algebra and Feynman Representation}
\label{chap:integral-genealogica}

Feynman's formulation assigns an amplitude to each possible history and sums those amplitudes to obtain the propagator. In HMT, the word \emph{history} has content prior to amplitude: a route preserves position, sheet, orientation, phase, carry, boundary, memory, and survival condition. The sum over paths is not added to that system from outside. It appears upon linearly representing its category of routes \cite{Feynman1948,FeynmanHibbs1965}.

The development has three linked mathematical levels. The expansion of a
matrix power as a finite sum of paths has already been proved in Dimensional
Mechanics. The ledger \(27\times27\) provides a unitary realization
with shift, mixer, and phase, diagonalizable by Fourier transform. The family of finite
sums subsequently admits a cylindrical limit when the measure, the action, and the
refinement satisfy the hypotheses of the convergence theorem. The
combinatorial identity, its spectral realization, and its functional extension
are thereby brought together.

\section{The enriched category of paths}

Let \(X\) be a finite set of observable states, and let \(\mathsf P_{\TPK}\) be the free category generated by the admissible TPK transitions. Denote \(\widetilde X=\operatorname{Ob}(\mathsf P_{\TPK})\). Its objects are enriched boundary states. A morphism \(\gamma:\widetilde x\to\widetilde y\) is a composable sequence of enriched transitions. Composition is concatenation, and the identity is the empty path.

The visible projection
\[
 q_{\rm vis}:\mathsf P_{\TPK}\longrightarrow \mathsf P_X
\]
forgets part of the fiber. Two different morphisms may have the same
observable word. This lack of faithfulness is essential: a sum that aggregated too
early over \(\mathsf P_X\) would lose precisely the information capable of
distinguishing incompatible futures.

\begin{definition}[Genealogical weight]
Let \(\mathcal B\) be a real normed algebra, not necessarily commutative, and
let \(B(\mathcal B,\cdot,1)\) be the one-object category whose endomorphism
monoid is \((\mathcal B,\cdot,1)\). A genealogical weight is a functor
\[
 W:\mathsf P_{\TPK}\longrightarrow B(\mathcal B,\cdot,1)
\]
which satisfies
\[
 W(\id_x)=1_{\mathcal B},\qquad
 W(\gamma_2\circ\gamma_1)=W(\gamma_2)W(\gamma_1).
\]
The order of the product follows the causal order of the route.
\end{definition}

The weight may contain a fibre transport, a positive penalty, an
action, a ledger character, or several of those factors. The definition does not
conflate those functions: each component retains its domain and composition
rule.

\section{The exact finite sum}

Let \(E_x\) be the linear fiber over \(x\in X\) and
\(\mathcal H_X=\bigoplus_{x\in X}E_x\). Let
\(U=(U_{yx})_{x,y\in X}\) be the global block operator, with
\(U_{yx}:E_x\to E_y\). We interpret \(U_{yx}\) as elementary transport
from \(x\) to \(y\). For a visible path
\(\gamma=(x_0,x_1,\ldots,x_R)\), we define
\[
 W_U(\gamma)=U_{x_Rx_{R-1}}\cdots U_{x_1x_0}.
\]

\begin{theorem}[Finite Path Expansion]
\label{thm:suma-finita-caminos}
For every \(R\geq 1\),
\begin{equation}
 (U^R)_{fi}
 =\sum_{\substack{\gamma:i\to f\\|\gamma|=R}}W_U(\gamma).
 \label{eq:suma-finita-caminos}
\end{equation}
If the global operator \(U:\mathcal H_X\to\mathcal H_X\) is orthogonal ---or
unitary after choosing a complex chart---, the
evolution \(\psi_R=U^R\psi_0\) preserves the norm.
\end{theorem}

\begin{proof}
The entry \((f,i)\) of the product \(U^R\) is
\[
 \sum_{x_1,\ldots,x_{R-1}}
 U_{fx_{R-1}}U_{x_{R-1}x_{R-2}}\cdots U_{x_1i}.
\]
Each choice of intermediate indices determines a unique path of length
\(R\), and each path determines a unique ordered monomial. Norm
conservation is an immediate consequence of \((U^R)^{\mathsf T}U^R=I\), or of its
unitary version.
\end{proof}

The formula \eqref{eq:suma-finita-caminos} is the combinatorial identity
underlying the sum over histories. Its HMT content appears when each
visible path is refined by the fibers of \(q_{\rm vis}\):
\begin{equation}
 \mathcal K_R(\widetilde f,\widetilde i)
 =\sum_{\substack{\gamma\in
        \mathsf P_{\TPK}(\widetilde i,\widetilde f)\\|\gamma|=R}}
 W(\gamma).
 \label{eq:kernel-genealogico-finito}
\end{equation}
In that sum, two published routes with the same endpoints remain distinct
terms whenever their genealogies differ.

\section{Unitary realization on the genealogical record \(27\times27\)}

The historical source develops a quantum walk on the torus
\[
 \Lambda_{27}=(\ZZ/27\ZZ)^2
\]
with four cardinal orientations \(\mathcal D=\{N,S,E,W\}\). Its real phase space is
\[
 \mathcal H_{27}^{\RR}
 =\ell^2(\Lambda_{27})\otimes\RR^4\otimes E,
\]
where \((E,J_E)\) internally realizes the complex structure. The direction
mixer is the Walsh--Hadamard matrix
\[
 C_4=\frac12
 \begin{pmatrix}
 1&1&1&1\\
 1&1&-1&-1\\
 1&-1&1&-1\\
 1&-1&-1&1
 \end{pmatrix},
 \qquad C_4^{\mathsf T}C_4=I_4.
\]
An oriented elementary action \(s_d\in\mathcal L_S\) defines the real
phase gate
\[
 D_J=\operatorname{diag}
 \left(
  e^{-J_Es_N/\hbar_{\rm int}},
  e^{-J_Es_S/\hbar_{\rm int}},
  e^{-J_Es_E/\hbar_{\rm int}},
  e^{-J_Es_W/\hbar_{\rm int}}
 \right).
\]
The conditioned shift transports each component in its direction:
\begin{align*}
 (S\psi)(x,y,N)&=\psi(x,y-1,N),&
 (S\psi)(x,y,S)&=\psi(x,y+1,S),\\
 (S\psi)(x,y,E)&=\psi(x-1,y,E),&
 (S\psi)(x,y,W)&=\psi(x+1,y,W),
\end{align*}
with coordinates modulo \(27\). The elementary step and the dodecaphase
composition are
\begin{equation}
 U_{\rm micro}=S D_J C_4,
 \qquad
 U_{12}=U_{\rm micro}^{12}.
 \label{eq:unitario-ledger-27}
\end{equation}

\begin{proposition}[Unitarity and exact sum in the ledger]
The operator \(U_{\rm micro}\) is orthogonal on the realification and unitary in
the complex chart. For every \(n\), each entry of \(U_{\rm micro}^n\) is
the ordered sum of the weights of all paths of \(n\) steps joining
its endpoint states.
\end{proposition}

\begin{proof}
The shift is a permutation, \(C_4\) is orthogonal, and each block
\(e^{-J_Es_d/\hbar_{\rm int}}\) is a rotation. Their product is orthogonal.
The second assertion is the
\cref{thm:suma-finita-caminos} applied to the block matrix
\(U_{\rm micro}\).
\end{proof}

Translational symmetry reduces the spectral problem to blocks of size
four. For
\[
 k_x=\frac{2\pi m}{27},\qquad
 k_y=\frac{2\pi n}{27},
 \qquad 0\le m,n<27,
\]
the shift diagonalizes as
\[
 S(k_x,k_y)=\operatorname{diag}
 \bigl(e^{J_Ek_y},e^{-J_Ek_y},e^{J_Ek_x},e^{-J_Ek_x}\bigr),
\]
and
\begin{equation}
 U_{\rm micro}(k_x,k_y)=S(k_x,k_y)D_JC_4,
 \qquad
 U_{12}(k_x,k_y)=U_{\rm micro}(k_x,k_y)^{12}.
 \label{eq:bloque-fourier-ledger-27}
\end{equation}
The global spectrum is the union of the spectra of the \(27^2\) blocks.
The expression combines two representations of the same dynamics: the matrix
power enumerates paths, and the Fourier transform publishes
quasienergies.

In each unit block a spectral logarithm can be chosen and defined
\[
 H_{\rm eff}(k_x,k_y)
 =\frac{\hbar_{\rm int}}{J_E\,\Delta t}
   \log U_{12}(k_x,k_y),
\]
understood through functional calculus in the complex chart or its realification.
The choice of branch fixes the quasienergy zone. The effective Hamiltonian
thus arises from the propagator already constructed; its exponential recovers
\(U_{12}\) in each mode.

\section{Real action character}

Let \(\mathcal A_*:\mathsf P_{\TPK}\to\RR_{\geq0}\) be the additive action cocycle of the
HMT--MD source. The gauge is
\[
 \lambda_*
 =1-\frac5{468}-\frac3{11}\Delta_4
 =0.9892859130191415\ldots\in(0,1).
\]
The local clause \(\mathsf H_{\rm loc}\) assigns unit cost to a turn,
cost \(\lambda_*\) to a change of sheet, and adds both when they occur on the
same edge. Under this clause,
\[
 \mathcal A_*(\gamma)
 =\#\operatorname{turns}(\gamma)
  +\lambda_*\#\operatorname{sheet\mbox{-}switches}(\gamma),
\qquad
 \mathcal A_*(\gamma_2\circ\gamma_1)
 =\mathcal A_*(\gamma_2)+\mathcal A_*(\gamma_1).
\]
The associated dimensional action is
\[
 S_{\rm int}(\gamma)=\hbar_{\rm int}\mathcal A_*(\gamma).
\]

In the real phase module \((E,J_E)\) of the preceding chapter, with \(E\) a
finite-dimensional real Euclidean space, define
\begin{equation}
 \Phi_J(\gamma)
 :=\exp\!\left(-J_E\frac{S_{\rm int}(\gamma)}{\hbar_{\rm int}}\right)
 =\cos\mathcal A_*(\gamma)\,I
  -\sin\mathcal A_*(\gamma)\,J_E.
 \label{eq:caracter-real-accion}
\end{equation}

\begin{proposition}[Multiplicativity of the internal phase]
\label{prop:multiplicatividad-fase-interna}
The map \(\Phi_J\) is a functor from \(\mathsf P_{\TPK}\) to the orthogonal
group generated by \(J_E\):
\[
 \Phi_J(\id_x)=I,
 \qquad
 \Phi_J(\gamma_2\circ\gamma_1)
 =\Phi_J(\gamma_2)\Phi_J(\gamma_1).
\]
\end{proposition}

\begin{proof}
The first equality uses \(\mathcal A_*(\id_x)=0\). For the second, the
additivity of \(\mathcal A_*\) and the commutation of multiples of the same
endomorphism \(J_E\) give
\[
 e^{-J_E(\mathcal A_*(\gamma_2)+\mathcal A_*(\gamma_1))}
 =e^{-J_E\mathcal A_*(\gamma_2)}
  e^{-J_E\mathcal A_*(\gamma_1)}.
\]
\end{proof}

Now let \(T(e)\in\End(E)\) be the bounded fiber transport of an edge, and
suppose
that it preserves \(J_E\). The weight
\[
 W_J(\gamma)
 =\prod_{e\in\gamma}^{\longleftarrow}
   T(e)\exp\!\left(-J_E\frac{s(e)}{\hbar_{\rm int}}\right)
\]
realizes, entirely on \(E\), the Hamiltonian convention \(\exp(-iS_H[\gamma]/\hbar)\). The usual Feynman Lagrangian convention \(\exp(+iS_L[\gamma]/\hbar)\) is obtained by taking \(S_L=-S_H\) or, equivalently, by introducing an orientation \(\varepsilon\in\{-1,+1\}\) in \(\exp(\varepsilon J_ES/\hbar_{\rm int})\). The scalar symbol \(i\) is the complex chart of the algebraic operation performed by \(J_E\). The positive cost \(\mathcal A_*\) becomes physical action upon application of the corresponding realization map.

\section{Composition of propagators over concatenated paths and amplitude law}

The genealogy preserves the composition law that makes the kernel dynamic.

\begin{proposition}[Convolution and Temporal Truncation]
\label{prop:chapman-kolmogorov-genealogico}
Suppose that \(\widetilde X\) is finite; more generally, suppose that the following family of
products is absolutely summable. If every route of length
\(R+S\) has a unique cut after \(R\) steps, then
\begin{equation}
 \mathcal K_{R+S}(\widetilde f,\widetilde i)
 =\sum_{\widetilde x\in\widetilde X}
   \mathcal K_S(\widetilde f,\widetilde x)
   \mathcal K_R(\widetilde x,\widetilde i).
 \label{eq:convolucion-kernel}
\end{equation}
The identity remains valid with noncommutative weights if the causal order of the factors is preserved.
\end{proposition}

\begin{proof}
We partition the set of routes by the enriched state
\(\widetilde x\) reached at the cut.
Functorial factorization of the weight transforms the sum over complete routes
into the product of the two partial sums, with the subsequent segment on the
left.
\end{proof}

The cut on the visible state would be valid only if weight and composition
descended faithfully through \(q_{\rm vis}\). In the presence of memory, two
segments with the same visible coordinate may belong to incompatible fibers; hence
convolution is performed over \(\widetilde X\).

This equation distinguishes the sum over paths from a collection of numerical
coincidences. The propagator has a semigroup, endpoints, composition, and a
refinement law.

\section{Cylindrical limit of the finite system and construction of the corresponding integral}

A functional integral requires control of the path space and its
refinement. For enriched endpoints \(\widetilde i,\widetilde f\), let
\(\mathcal P_n(\widetilde i,\widetilde f)\) be the finite sets of admissible paths
at resolution \(n\), with surjective restrictions, and define
\[
 \Omega_{fi}^{\rm path}
 :=\varprojlim_n\mathcal P_n(\widetilde i,\widetilde f).
\]
It is the compact space of path genealogies with those boundary data.
Let \(\{\Pi_n\}_{n\geq0}\) be a sequence of cylindrical partitions of
\(\Omega_{fi}^{\rm path}\), with \(\Pi_{n+1}\) refining \(\Pi_n\). For each
cylinder \(C\in\Pi_n\), let us choose a representative
\(\gamma_C\), a projective weight \(\mu_n(C)\), and an approximate action
\(S_n(C)\).

We define the simple operator
\begin{equation}
 \mathcal I_n(F)
 :=\sum_{C\in\Pi_n}
   \mu_n(C)
   F(\gamma_C)
   \exp\!\left(-J_E\frac{S_n(C)}{\hbar_{\rm int}}\right).
 \label{eq:integral-cilindrica-n}
\end{equation}

\begin{theorem}[Genealogical Limit of Cylindrical Integrals]
\label{thm:integral-cilindrica-limite}
Suppose that:
\begin{enumerate}
 \item the cylindrical masses are projectively compatible;
 \item \(F:\Omega_{fi}^{\rm path}\to\End(E)\) is continuous and bounded;
 \item there exists a continuous function
       \(S:\Omega_{fi}^{\rm path}\to\mathcal L_S\) such that
       \[
        \max_{C\in\Pi_n}\sup_{\gamma\in C}
        \left\|
          \frac{S_n(C)-S(\gamma)}{\hbar_{\rm int}}
        \right\|\longrightarrow0;
       \]
 \item the mesh of the partitions,
       \(\operatorname{mesh}(\Pi_n)
        =\max_{C\in\Pi_n}\diam_{d_\vartheta}C\),
       tends to zero.
\end{enumerate}
Since \(E\) is finite-dimensional, \(\End(E)\), endowed with the operator
norm, is a Banach space; the following integral is therefore a well-defined
Bochner integral \cite{Bochner1955}.
Then \(\mathcal I_n(F)\) converges in operator norm, and its limit does not
depend on the representatives. We write
\begin{equation}
 \int_{\Omega_{fi}^{\rm path}}
 F(\gamma)
 \exp\!\left(-J_E\frac{S(\gamma)}{\hbar_{\rm int}}\right)
 \dd\mu(\gamma)
 :=\lim_{n\to\infty}\mathcal I_n(F).
 \label{eq:integral-genealogica}
\end{equation}
\end{theorem}

\begin{proof}
Projective compatibility constructs the cylinder measure \(\mu\). The
function
\[
 G(\gamma)=F(\gamma)
 \exp\!\left(-J_E S(\gamma)/\hbar_{\rm int}\right)
\]
is continuous on the compact set \(\Omega_{fi}^{\rm path}\), hence uniformly
continuous and bounded. Let
\[
 G_n(C)=F(\gamma_C)
 \exp\!\left(-J_E S_n(C)/\hbar_{\rm int}\right).
\]
Since \(J_E\) is an orthogonal structure,
\[
 \|e^{-J_Ea}-e^{-J_Eb}\|\leq |a-b|.
\]
The third hypothesis implies that
\(G_n(C)-G(\gamma_C)\) tends uniformly to zero. The fourth hypothesis and the
uniform continuity of \(G\) make the sums with \(G(\gamma_C)\)
convergent Riemann sums independent of the representatives. The
limit coincides with the Bochner integral of \(G\) with respect to \(\mu\).
More precisely, if
\[
 \epsilon_n
 =\max_{C\in\Pi_n}\sup_{\gamma\in C}
   \left\|
    \frac{S_n(C)-S(\gamma)}{\hbar_{\rm int}}
   \right\|,
\]
then
\[
 \left\|\mathcal I_n(F)-\int G\,\dd\mu\right\|
 \leq
 \omega_F\!\bigl(\operatorname{mesh}(\Pi_n)\bigr)
 +\|F\|_\infty\epsilon_n,
\]
because the integral decomposes over the cylinders and the total mass is one.
Both terms tend to zero by hypothesis.
\end{proof}

\begin{corollary}[Convergence of a family of discrete kernels]
\label{cor:convergencia-kernels-discretos}
In addition to the hypotheses of
\cref{thm:integral-cilindrica-limite}, suppose that there exists a cofinal sequence
of depths \(R_n\) and simple functions \(F_{fi,n}\) such that
\[
 \mathcal K_{R_n}(\widetilde f,\widetilde i)
 =\mathcal I_n(F_{fi,n}),
 \qquad
 F_{fi,n}\longrightarrow F_{fi}
\]
uniformly. Then the kernels
\(\mathcal K_{R_n}(\widetilde f,\widetilde i)\) converge to the genealogical integral
of \(F_{fi}\).
\end{corollary}

\begin{proof}
The difference between \(\mathcal I_n(F_{fi,n})\) and
\(\mathcal I_n(F_{fi})\) is bounded by
\(\|F_{fi,n}-F_{fi}\|_\infty\), because \(\mu_n\) has total mass one and the
phase is orthogonal. The second term converges by
\cref{thm:integral-cilindrica-limite}.
\end{proof}

The theorem produces a rigorous integral over the compact genealogical space. The corollary makes explicit the bridge that identifies each matrix kernel with a normalized cylindrical sum. The real-time Feynman oscillatory integral appears when a physical representation and a refinement whose discrete action approximates the action of the system are chosen. That identification requires proof of convergence of the corresponding propagator; it is not granted by the mere existence of \(\mu\).

\section{Integral visible and integral conditioned by memory}

A real path reader \(\Val_{\rm path}:\Omega_{fi}^{\rm path}\to K\subset\RR\) induces the visible measure
\(\nu=(\Val_{\rm path})_*\mu\). For a function that factors through the value,
\(F=\widehat F\circ\Val_{\rm path}\),
\[
 \int_{\Omega_{fi}^{\rm path}}F\,\dd\mu
 =\int_K\widehat F(x)\,\dd\nu(x).
\]
The phase or transport, however, may fail to factor through
\(\Val_{\rm path}\). The disintegration of \(\mu\) over the value fibers
expresses the difference:
\begin{equation}
 \int_{\Omega_{fi}^{\rm path}}G(\gamma)\,\dd\mu(\gamma)
 =\int_K
   \left[
    \int_{\Val_{\rm path}^{-1}(x)}G(\gamma)\,\dd\mu_x(\gamma)
   \right]\dd\nu(x).
 \label{eq:desintegracion-caminos}
\end{equation}
The inner integral preserves the histories that publish the same coordinate.
This is the mathematical form of integration with respect to values without reducing the
number to its digit.

\section{Action, interference, and persistence}

The genealogical sum simultaneously maintains three operations:
\[
 \text{route admissibility}
 \quad\longrightarrow\quad
 \text{weight/action}
 \quad\longrightarrow\quad
 \text{interference in the representation}.
\]
The first belongs to TPK and its survival law. The second belongs
to the cocycle and dimensional transport. The third belongs to the phase
module and the physical reader. No observed magnitude retrospectively
selects the route that produces it.

\begin{resultbox}
The finite Feynman sum is obtained as an exact linear representation of the
enriched category of routes. The complex structure can be realized by a real
endomorphism \(J_E^2=-I\). Under compatibility of measures, contraction of
cylinders, and uniform convergence of actions, the cylindrical sums define
a genealogical integral. The finite kernels converge to it when they
additionally satisfy the bridge of the
\cref{cor:convergencia-kernels-discretos}.
\end{resultbox}

\begin{statusbox}
The finite matrix expansion, the action cocycle, and the real realizations
of the complex structure are recovered results. Their assembly in
\eqref{eq:caracter-real-accion} and \eqref{eq:kernel-genealogico-finito} is an
exact new formalization. The theorem \ref{thm:integral-cilindrica-limite} is
exact under its premises. Equivalence with a concrete continuous quantum model
requires verifying those premises for the Hamiltonian and the action of
that model.
\end{statusbox}

The premises of the theorem identify exactly the structure that passes from the finite ledger to the function space: projective compatibility of the weights, contraction of the mesh, uniform convergence of the action, and correspondence between kernels and cylindrical sums. In the realization \eqref{eq:unitario-ledger-27}, composition and unitarity are algebraic identities; as the ledger is refined, the remaining premises determine convergence to the chosen continuum propagator.
